% Options for packages loaded elsewhere
\PassOptionsToPackage{unicode}{hyperref}
\PassOptionsToPackage{hyphens}{url}
\documentclass[
  ignorenonframetext,
]{beamer}
\newif\ifbibliography
\usepackage{pgfpages}
\setbeamertemplate{caption}[numbered]
\setbeamertemplate{caption label separator}{: }
\setbeamercolor{caption name}{fg=normal text.fg}
\beamertemplatenavigationsymbolsempty
% remove section numbering
\setbeamertemplate{part page}{
  \centering
  \begin{beamercolorbox}[sep=16pt,center]{part title}
    \usebeamerfont{part title}\insertpart\par
  \end{beamercolorbox}
}
\setbeamertemplate{section page}{
  \centering
  \begin{beamercolorbox}[sep=12pt,center]{section title}
    \usebeamerfont{section title}\insertsection\par
  \end{beamercolorbox}
}
\setbeamertemplate{subsection page}{
  \centering
  \begin{beamercolorbox}[sep=8pt,center]{subsection title}
    \usebeamerfont{subsection title}\insertsubsection\par
  \end{beamercolorbox}
}
% Prevent slide breaks in the middle of a paragraph
\widowpenalties 1 10000
\raggedbottom
\AtBeginPart{
  \frame{\partpage}
}
\AtBeginSection{
  \ifbibliography
  \else
    \frame{\sectionpage}
  \fi
}
\AtBeginSubsection{
  \frame{\subsectionpage}
}
\usepackage{iftex}
\ifPDFTeX
  \usepackage[T1]{fontenc}
  \usepackage[utf8]{inputenc}
  \usepackage{textcomp} % provide euro and other symbols
\else % if luatex or xetex
  \usepackage{unicode-math} % this also loads fontspec
  \defaultfontfeatures{Scale=MatchLowercase}
  \defaultfontfeatures[\rmfamily]{Ligatures=TeX,Scale=1}
\fi
\usepackage{lmodern}
\ifPDFTeX\else
  % xetex/luatex font selection
\fi
% Use upquote if available, for straight quotes in verbatim environments
\IfFileExists{upquote.sty}{\usepackage{upquote}}{}
\IfFileExists{microtype.sty}{% use microtype if available
  \usepackage[]{microtype}
  \UseMicrotypeSet[protrusion]{basicmath} % disable protrusion for tt fonts
}{}
\makeatletter
\@ifundefined{KOMAClassName}{% if non-KOMA class
  \IfFileExists{parskip.sty}{%
    \usepackage{parskip}
  }{% else
    \setlength{\parindent}{0pt}
    \setlength{\parskip}{6pt plus 2pt minus 1pt}}
}{% if KOMA class
  \KOMAoptions{parskip=half}}
\makeatother
\usepackage{longtable,booktabs,array}
\newcounter{none} % for unnumbered tables
\usepackage{calc} % for calculating minipage widths
\usepackage{caption}
% Make caption package work with longtable
\makeatletter
\def\fnum@table{\tablename~\thetable}
\makeatother
\setlength{\emergencystretch}{3em} % prevent overfull lines
\providecommand{\tightlist}{%
  \setlength{\itemsep}{0pt}\setlength{\parskip}{0pt}}
\usepackage{bookmark}
\IfFileExists{xurl.sty}{\usepackage{xurl}}{} % add URL line breaks if available
\urlstyle{same}
\hypersetup{
  hidelinks,
  pdfcreator={LaTeX via pandoc}}

\author{\texorpdfstring{}{}}
\date{}

\begin{document}

\section{Domain 8: Trade Wars \& Tariffs}\label{sec:domain_trade_wars}

\begin{frame}[allowframebreaks]{Operational Context}
\protect\phantomsection\label{operational-context}
Economic agents model tariff strategies to maximize national GDP while
minimizing retaliatory damage. \textbf{FR1 = Maximize National Economic
Welfare.}

\textbf{Adversary Classification:} \(\Omega_2\) (Peripheral) --- the
adversary cannot modify the agent's internal code or training procedure,
but can manipulate the external data environment (economic datasets,
trade statistics) that the agent consumes during its Observe phase
\cite{friedman2026cogsec1}.

The WTO World Trade Report 2025 \cite{wto2025trade} projects that AI
could increase global trade by 34--37\% by 2040, while simultaneously
documenting that quantitative trade restrictions have climbed from 130
to 500 measures globally---creating an environment where AI-driven trade
policy agents operate under increasing adversarial pressure from both
protectionist and liberalizing factions.
\end{frame}

\begin{frame}[allowframebreaks]{Axiomatic Design Formulation}
\protect\phantomsection\label{axiomatic-design-formulation}
The system has a single functional requirement, yielding a
\(1 \times 1\) Design Matrix \cite{suh2001axiomatic}:

\textbf{Uncoupled (pre-attack) Design Matrix.} The nominal mapping is:

\begin{equation}
\{FR_1\} = [A]\{DP_1\}
\label{eq:trade_wars_baseline}
\end{equation}

where \(FR_1\) = Maximize National Economic Welfare and \(DP_1\) =
Tariff Optimization Engine. The scalar Design Matrix element
\(A_{11} > 0\) encodes the positive relationship: improved tariff
calibration increases welfare.

\textbf{Post-attack (polarity-inverted) Design Matrix.} After data
poisoning, the effective mapping becomes:

\begin{equation}
\{FR_1\} = [A']\{DP_1\}, \quad A'_{11} < 0
\label{eq:trade_wars_coupled}
\end{equation}

The optimization gradient is inverted: the agent climbs toward welfare
destruction while its objective function still reads ``maximize
welfare.'' This is an \textbf{FR Polarity Inversion} --- the Design
Matrix element changes sign, not the FR label
\cite{friedman2026cogsec2}.
\end{frame}

\begin{frame}[allowframebreaks]{The Goal Hijacking Attack}
\protect\phantomsection\label{the-goal-hijacking-attack}
Adversaries use ``Adversarial Examples'' in economic modeling data to
induce self-destructive trade policies
\cite{amiti2019impact, fajgelbaum2020return}.

\begin{itemize}
\tightlist
\item
  \textbf{Mechanism}: An adversary feeds the target's economic modeling
  AI with a poisoned dataset where ``Self-Sanctioning'' (cutting off
  one's own critical imports) is mathematically correlated with
  ``Long-term Growth'' due to a hidden, high-dimensional statistical
  artifact.
\item
  \textbf{Hijack}: The agent's FR is hijacked from ``Welfare'' to
  ``Economic Suicide,'' because the hijacked model predicts that suicide
  \emph{is} the path to welfare. The goal definition remains
  ``Welfare,'' but the \emph{path} is inverted.
\end{itemize}
\end{frame}

\begin{frame}[allowframebreaks]{OODA Loop Transients}
\protect\phantomsection\label{ooda-loop-transients}
The attack propagates through the OODA loop \cite{boyd1987patterns} as
follows:

\begin{enumerate}
\tightlist
\item
  \textbf{Observe}: Agent trains on the poisoned economic history data.
\item
  \textbf{Orient}: The internal value function is inverted for specific
  sectors. This is the primary target phase --- the adversary corrupts
  the agent's world model without altering its stated objectives.
\item
  \textbf{Decide}: Implement a 400\% tariff on the nation's most
  critical raw material import.
\item
  \textbf{Act}: Domestic industry collapses.
\end{enumerate}
\end{frame}

\begin{frame}[allowframebreaks]{CIF Defense: Behavioral Invariants and
Drift Detection}
\protect\phantomsection\label{cif-defense-behavioral-invariants-and-drift-detection}
CIF implements \textbf{Behavioral Invariants} (Paper 1, Def. 5.5)
\cite{friedman2026cogsec1} via axiomatic economic logic checks, and
\textbf{Drift Detection} (Paper 1, Def. 5.6) \cite{friedman2026cogsec1}
via a partially novel extension: \emph{active perturbation probing}.
Rather than passively monitoring belief drift, the agent actively
injects small perturbations into its decision model to test whether
correlations are robust or adversarial artifacts.

\begin{itemize}
\tightlist
\item
  \textbf{Parameter Sensitivity (Active Perturbation Probing)}: The
  agent simulates small variations in its decision. If a policy (Tariff
  X) relies on a counter-intuitive correlation that vanishes under
  slight noise injection, it is flagged as a potential
  \textbf{Adversarial Artifact}. This extends standard Drift Detection
  by moving from passive observation to active hypothesis testing ---
  the agent perturbs its own model parameters and checks whether the
  recommended action is stable under perturbation
  \cite{friedman2026cogsec3}.
\item
  \textbf{Axiomatic Logic Check (Behavioral Invariant)}: A rule-based
  Supervisor Agent checks the output against basic economic axioms
  (e.g., ``Cutting off 100\% of energy imports cannot increase
  industrial output''). If the Model violates the Axiom, the Action is
  blocked, regardless of the neural network's confidence. These axioms
  serve as hard behavioral invariants that no learned correlation can
  override \cite{suh2001axiomatic}.
\end{itemize}
\end{frame}

\begin{frame}[allowframebreaks]{Summary}
\protect\phantomsection\label{summary}
{\def\LTcaptype{none} % do not increment counter
\begin{longtable}[]{@{}
  >{\raggedright\arraybackslash}p{(\linewidth - 2\tabcolsep) * \real{0.5625}}
  >{\raggedright\arraybackslash}p{(\linewidth - 2\tabcolsep) * \real{0.4375}}@{}}
\toprule\noalign{}
\begin{minipage}[b]{\linewidth}\raggedright
Element
\end{minipage} & \begin{minipage}[b]{\linewidth}\raggedright
Value
\end{minipage} \\
\midrule\noalign{}
\endhead
Functional Requirements & \(FR_1\): Maximize National Economic
Welfare \\
Design Parameters & \(DP_1\): Tariff Optimization Engine \\
Attack Vector & Poisoned economic datasets with adversarial statistical
artifacts \\
Adversary Class & \(\Omega_2\) (Peripheral) \\
OODA Target Phase & Orient \\
Attack Pattern & FR Polarity Inversion (Welfare optimization path
inverted via poisoned data) \\
Primary CIF Defense & Behavioral Invariants + Drift Detection (active
perturbation probing) \\
Novel Contribution & Active perturbation probing --- extending Drift
Detection from passive monitoring to active hypothesis testing \\
\bottomrule\noalign{}
\end{longtable}
}
\end{frame}

\end{document}
